\section{Deadline-aware and Mobile CPU Scheduling}
\label{sec:related-deadline-aware-and-mobile-cpu-scheduling}

전통적인 Real-time Scheduling 연구는 명시적인 Task Timing Constraint를 기반으로 실행 순서를 결정한다. Liu와 Layland는 주기적 Task Model에서 Earliest Deadline First (EDF)의 Scheduling 특성을 분석하였으며~\cite{liu1973}, Constant Bandwidth Server (CBS)는 Multimedia Task와 Hard Real-time Task의 CPU 자원 공유를 지원한다~\cite{abeni1998}. Linux의 SCHED\_DEADLINE은 EDF와 CBS를 기반으로 Runtime, Deadline 및 Period를 지정하는 Interface를 제공한다~\cite{scheddeadline}. 또한 Priority Inheritance는 높은 Priority를 가진 Task가 Lock Dependency로 인해 지연되는 문제를 완화한다~\cite{sha1990}. 이러한 접근은 명시적인 시간 제약이나 Synchronization 정보를 활용하는 반면, RL-SchedExt는 Application Request의 Deadline을 Runtime에 직접 제공받지 않는 일반 Task의 실행 순서를 조정한다. 따라서 본 연구는 Hard Real-time Guarantee를 제공하는 EDF나 CBS의 대체가 아니라, 일반 Android Workload의 Deadline Compliance를 개선하는 접근이다.

Android의 Interactivity를 개선하기 위한 연구도 수행되었다. Huh 등은 Android Framework의 정보를 활용하여 Interactive Task를 식별하고 Virtual Time 기반의 CFS를 개선하였다~\cite{huh2015}. Zheng 등은 실제 Mobile System의 느린 UI 응답 원인을 분석하고 Process Management를 개선하는 접근을 제안하였다~\cite{zheng2024}. 한편 CFS와 FreeBSD ULE의 비교 연구는 Scheduler의 설계 선택에 따라 Application Performance가 달라질 수 있음을 보여준다~\cite{battle}. RL-SchedExt는 Application Framework가 전달하는 명시적인 Request Deadline이나 Task Category 대신 Kernel에서 관측 가능한 Runtime State를 이용하고, 경쟁하는 Runnable Task Set의 Context에 따라 Learned Vruntime Adjustment를 계산한다.

\section{Extensible and Latency-sensitive Scheduling}
\label{sec:related-extensible-and-latency-sensitive-scheduling}

ghOSt는 Linux Scheduling Policy를 User Space에 위임하여 다양한 Scheduler의 개발과 배포를 지원한다~\cite{ghost}. Enoki는 Linux Kernel Scheduler의 개발 및 검증 과정을 개선하는 Framework를 제안하였다~\cite{enoki}. Linux sched\_ext는 eBPF Scheduling Program을 통해 사용자 정의 정책을 Kernel에서 실행할 수 있는 Interface를 제공한다~\cite{schedext}. 특히 scx\_lavd는 Task의 실행 및 Wakeup 특성 등으로부터 Latency Criticality를 추정하여 Virtual Deadline을 결정하는 sched\_ext 기반 Scheduler이다~\cite{lavd}. scx\_lavd와 RL-SchedExt는 실제 Application Deadline을 직접 입력으로 요구하지 않는다는 공통점이 있다. 차이는 RL-SchedExt가 Simulator에서 Deadline-oriented Reward를 최적화한 Nonlinear Scheduling Policy를 이용하여 CFS 기반 Vruntime에 대한 Adjustment를 계산한다는 점이다. 본 논문의 실험에는 scx\_lavd를 포함하지 않았으므로 두 정책 사이의 성능 우위는 주장하지 않는다.

Datacenter 환경에서도 Tail Latency를 개선하기 위한 다양한 Scheduling 연구가 제안되었다. Shinjuku는 짧은 시간 단위의 Preemption을 통해 실행 시간이 다양한 Request의 Tail Latency를 개선하며~\cite{shinjuku}, Shenango는 빠른 Core Reallocation을 통해 Latency-sensitive Application과 Batch Application 사이의 CPU 자원 활용을 개선한다~\cite{shenango}. Caladan은 미세한 시간 단위로 자원 간 Interference를 감지하고 Task Placement를 조정한다~\cite{caladan}. 이들 연구와 달리 RL-SchedExt의 현재 평가 범위는 Android의 단일 CPU 내 Task Ordering이며, Core Allocation이나 Multi-core Load Balancing을 최적화하지 않는다.

\section{Learning-based Scheduling and In-Kernel Inference}
\label{sec:related-learning-based-scheduling-and-in-kernel-inference}

DeepRM은 RL을 이용하여 다차원 자원 요구량을 가진 Job의 Resource Allocation Policy를 학습하며~\cite{deeprm}, Decima는 Job의 Dependency Graph를 반영하여 Data Processing Cluster의 Scheduling Policy를 학습한다~\cite{decima}. 이러한 연구는 Workload 특성에 따라 Scheduling Policy를 학습할 수 있음을 보여준다. RL-SchedExt는 Cluster-level Job Scheduling과 달리 빈번하게 발생하는 Kernel Scheduling Decision에서 개별 Runnable Task의 실행 순서를 결정하며, Runtime에 Application-level Dependency Graph가 제공되는 것을 전제로 하지 않는다.

Linux 및 Android Scheduler에 직접 RL을 적용한 연구로 Learning EAS, STUN 및 MobiRL이 있다. Learning EAS는 Policy Gradient를 이용하여 CPU Frequency 설정과 Task Migration에 관련된 Scheduler Parameter를 조정하며~\cite{han2020}, STUN은 RL을 이용하여 주어진 Workload에 적합한 Linux Scheduler Parameter를 탐색한다~\cite{stun}. MobiRL은 Android의 UI Smoothness와 Power Consumption을 함께 최적화하기 위해 CPU/GPU Frequency를 조정한다~\cite{mobirl}. 반면 RL-SchedExt의 Action은 개별 Task의 Learned Vruntime Adjustment이며, Task Selection 자체를 변경한다. 따라서 본 연구는 Frequency Scaling이나 Scheduler Parameter Tuning과는 다른 Scheduling 제어 단위를 대상으로 한다.

LinnOS는 경량 Neural Network를 Operating System에 통합하여 개별 Storage I/O의 성능을 추론하는 접근을 제시하였다~\cite{linnos}. RL-SchedExt 역시 빈번한 Kernel Decision에서 Inference Cost를 고려해야 한다는 문제를 공유한다. 본 연구에서는 RL Teacher의 Policy를 Neural Student로 Distillation하고 INT8 Quantization을 적용하여 Native Kernel에서 Inference를 수행하며, eBPF Scheduling Program은 그 결과를 이용하여 Task Selection을 수행한다. 이러한 설계는 학습 기반 Task Ordering을 실제 Android Kernel에 적용하기 위한 것으로, 학습 및 Runtime의 분리 자체보다 Deadline 정보를 직접 관측할 수 없는 조건에서 Nonlinear Scheduling Policy를 이용한 Task Ordering의 효과를 검증하는 데 초점을 둔다.
